\chapter{Interactive Proofs, Arithmetization and the PCP Theorem}
\label{ch:pcp}

The barriers of Chapter~\ref{ch:barriers} say that relativizing and algebrizing proofs cannot resolve $\Pc$ versus $\NP$. The results that do not relativize, the ones that use the algebraic structure of computation rather than treating it as a black box, mostly come from one idea, arithmetization, and the theory built on it, interactive proofs and probabilistically checkable proofs, is the deepest thing known about $\NP$. This chapter gives that theory in outline and explains its two connections to the program: it clarifies what the barriers do and do not say about derandomization, and the PCP theorem is a statement about the robustness of $\NP$ witnesses whose bearing on level sets is suggestive but not established here.

\section{Interactive proofs}

An interactive proof, introduced by Goldwasser, Micali and Rackoff in 1985 and in a public coin form by Babai in the same year, replaces the static certificate of $\NP$ by a conversation: a polynomial time verifier with random coins exchanges messages with an all powerful prover, and accepts or rejects. The class $\mathrm{IP}$ consists of languages for which a prover can convince the verifier of true statements with probability at least $\tfrac23$ and no prover can convince it of false ones with probability more than $\tfrac13$. Graph non isomorphism is in $\mathrm{IP}$ by a two round protocol, which was the first evidence that interaction adds power, since it is not known to be in $\NP$.

Shamir's theorem of 1992, building on Lund, Fortnow, Karloff and Nisan, states that $\mathrm{IP} = \mathrm{PSPACE}$: interaction with randomness is exactly as powerful as polynomial space. The proof is by arithmetization. A quantified Boolean formula is converted to a polynomial identity over a finite field by replacing AND with multiplication, OR with the inclusion exclusion polynomial, and the quantifiers with sums and products over $\{0, 1\}$, with Shamir's degree reduction after each quantifier to keep the degree polynomial; the prover claims the value of the resulting sum, and the verifier checks it by the sum check protocol, asking the prover to send the partial sum as a polynomial in one variable at a time and substituting a random field element, so that a false claim is caught with probability at least one minus the degree over the field size at each round. The protocol does not relativize because the arithmetization of an oracle gate is not available.

\section{The PCP theorem}

A probabilistically checkable proof is a certificate that the verifier reads only in a few places. The PCP theorem of Arora, Safra, Lund, Motwani, Sudan and Szegedy in 1992, with a new proof by Dinur in 2007, states that every language in $\NP$ has a proof system in which the verifier uses $O(\log n)$ random bits, reads a constant number of bits of the proof, accepts every valid proof of a true statement, and rejects every purported proof of a false statement with probability at least one half. In symbols, $\NP = \mathrm{PCP}(\log n, 1)$.

The theorem is equivalent to a statement about approximation: there is a constant $\epsilon > 0$ such that it is $\NP$ hard to distinguish satisfiable $3$ SAT formulas from those in which no assignment satisfies more than a $1 - \epsilon$ fraction of clauses. This gap amplification is Dinur's route to the proof: she shows that a formula can be transformed in polynomial time so that its unsatisfiability, if any, becomes a constant fraction of unsatisfied clauses, by alternating a graph powering step that amplifies the gap and an alphabet reduction step that keeps the alphabet small. The constant has since been pushed to the optimal values for several problems by H{\aa}stad's tight inapproximability results and, conditionally on the unique games conjecture of Khot, for others.

\section{Arithmetization and derandomization}

Hardness versus randomness, by contrast, relativizes. The Nisan Wigderson generator of Chapter~\ref{ch:pseudo} evaluates a hard function on overlapping subsets of a seed, and the analysis that a distinguisher yields a small circuit for the hard function is a black box reduction. The same holds for the hardness amplification that turns worst case hardness into average case hardness, in Impagliazzo and Wigderson's derandomized XOR lemma and in the list decoding approach of Sudan, Trevisan and Vadhan: these encode the hard function as a low degree polynomial and use the error correcting properties of polynomial codes, but the argument treats the function as a black box, and relativized forms of it are used by Klivans and van Melkebeek to derandomize AM from circuit hardness relative to an oracle. Using polynomial codes does not by itself make a proof non relativizing. The non relativizing results tied to derandomization are those of Impagliazzo, Kabanets and Wigderson and of Kabanets and Impagliazzo, which go through interactive proof characterizations of the classes involved. What the algebrization barrier of Aaronson and Wigderson blocks is the step that supplies the hardness, namely proving the circuit lower bound that the generator needs.

For the program the consequence is that any proof of $\Pc = \RP$ along the known lines, a hitting set generator from a hard function, reduces to proving a circuit lower bound, and, by Aaronson and Wigderson~\cite{AaronsonWigderson}, even the weaker lower bound $\mathrm{NEXP} \not\subseteq \mathrm{P/poly}$ cannot be proved by algebrizing techniques alone. The vertical view adds nothing to this and claims nothing; it says what the hitting set must hit.

\section{PCP and the geometry of level sets}

The second connection is more suggestive than established. The PCP theorem says that $\NP$ witnesses can be encoded so that a wrong witness for a false statement is wrong almost everywhere: for an unsatisfiable instance a purported proof fails a constant fraction of the local checks. Translated to energies, the PCP encoding of an instance is an energy on proofs, the number of failed checks, and soundness says that if the instance is unsatisfiable then every state has energy at least $\epsilon$ times the number of checks. For a satisfiable instance soundness says nothing about the non valid proofs: a valid proof with one bit flipped fails only the checks reading that bit, so near solutions exist. The gap is therefore between instances and not between states of one instance. For an unsatisfiable instance the level sets $\{ S \le m \}$ for $m$ below $\epsilon \cdot (\text{number of checks})$ are empty, which is the correct answer (the count is zero) and not a failure of the walk.

This is a hardness of approximation statement, and for Monte Carlo it poses a question rather than answering one: could a reduction produce boards whose energy has a gap, and would such a gap be a failure of the walk at all? The queens are not PCP encoded; their staircase has a step at every energy and the ladder exists. We do not know whether the Gent Jefferson Nightingale reduction applied to a PCP encoded formula produces a board whose energy, the number of attacking pairs, has a gap: that energy is not the number of violated PCP checks, and we do not prove that the reduction preserves any gap. If a reduction did preserve a gap, the empty bottom levels on boards from unsatisfiable instances would be the correct output, and boards from satisfiable instances would have near solutions as usual, so a gap would not by itself be a failure of the walk. We do not know of a mechanism by which it could be one.

\section{What fails on hard boards?}

For the program it matters what operates on hard boards, because the diagnostics differ. Shattering shows up as coalescence in the genealogy at the bottom rung and as torpid mixing of the level set chain; the walk reaches the near solutions and cannot cross between components. A gap, if it occurred, would show up as an empty level, which on a board from an unsatisfiable instance is the right answer and not a failure. On a reduction board from an ordinary hard formula it is tempting to say that the near solutions exist because a formula with one unsatisfied clause is a near solution, but near solutions of a formula and low energy configurations of the board are different objects, and whether shattering operates is a question for the experiments and not something we derive. Shattering is consistent with $\NP \neq \RP$ and with the thesis, which says the hardness is in the staircase's widths; the PCP theorem does not by itself say that any energy of interest has a gap or small widths.
